% Source: physical pp. 132--133 (printed pp. 109--110). Coverage: Problems 1--4 and References 1--5 only, beginning at the Problems heading after Section 17.5 and ending before Chapter 18. Figures/tables/footnotes/dividers: none. Uncertainties: none.

\chapterexercisehook{17}

\chapterbackmatter{References}

\begin{enumerate}
  \item \phantomsection\label{ch17-ref-1}
  J.~Zinn-Justin, in \textit{Trends in Elementary Particle Theory International Summer Institute on Theoretical Physics in Bonn 1974} (Springer-Verlag, Berlin, 1975).

  \item \phantomsection\label{ch17-ref-2}
  Before the advent of BRST symmetry, proofs of the renormalizability of non-Abelian gauge theories were based directly on the Slavnov--Taylor identities for gauge transformations; see B.~W.~Lee and J.~Zinn-Justin, \textit{Phys. Rev.} D\textbf{5}, 3121, 3137 (1972); \textit{Phys. Rev.} D\textbf{7}, 1049 (1972); G.~'t~Hooft and M.~Veltman, \textit{Nucl. Phys.} B\textbf{50}, 318 (1972); B.~W.~Lee, \textit{Phys. Rev.} D\textbf{9}, 933 (1974). The original proof of renormalizability based on BRST symmetry was given by C.~Becchi, A.~Rouet, and R.~Stora, \textit{Commun. Math. Phys.} \textbf{42}, 127 (1975); in \textit{Renormalization Theory---Proceedings of the International School of Mathematical Physics at Erice, August 1975}, eds. G.~Velo and A.S.~Wightman (D.~Reidel, Dordrecht, 1976): pp.~269--97, 299--343. The proof given here follows the general outline of J.~Zinn-Justin, Ref.~\hyperref[ch17-ref-1]{1}; B.~W.~Lee, in \textit{Methods in Field Theory}, eds. R.~Balian and J.~Zinn-Justin (North-Holland, Amsterdam, 1976): pp.~79--139.

  \item \phantomsection\label{ch17-ref-3}
  The point of view and presentation here is based on that of J.~Gomis and S.~Weinberg, Kyoto-Texas preprint RIMS-1036/UTTG-18-95 (1995), to be published in \textit{Nuclear Physics B}. For earlier use of these methods, see B.~L.~Voronov and I.~V.~Tyutin, \textit{Theor. Math. Phys.} \textbf{50}, 218 (1982); \textbf{52}, 628 (1982); B.~L.~Voronov, P.~M.~Lavrov, and I.~V.~Tyutin, \textit{Sov. J. Nucl. Phys.} \textbf{36}, 292 (1982); P.~M.~Lavrov and I.~V.~Tyutin \textit{Sov. J. Nucl. Phys.} \textbf{41}, 1049 (1985); D.~Anselmi, \textit{Class. and Quant. Grav.} \textbf{11}, 2181 (1994); \textbf{12}, 319 (1995); M.~Harada, T.~Kugo, and K.~Yamawaki, \textit{Prog. Theor. Phys.} \textbf{91}, 801 (1994).

  \item \phantomsection\label{ch17-ref-4}
  G.~Barnich and M.~Henneaux, \textit{Phys. Rev. Lett.} \textbf{72}, 1588 (1994); G.~Barnich, F.~Brandt, and M.~Henneaux, \textit{Phys. Rev.} \textbf{51}, R143 (1995); \textit{Commun. Math. Phys.} \textbf{174}, 57, 93 (1995); \textit{Nucl. Phys.} B\textbf{455}, 357 (1995).

  \item \phantomsection\label{ch17-ref-5}
  The background field gauge was introduced by B.~S.~De Witt, \textit{Phys. Rev.} \textbf{162}, 1195, 1239 (1967). For the treatment of multi-loop effects, see G.~'t~Hooft, in \textit{Functional and Probabilistic Methods in Quantum Field Theory: Proceedings of the 12th Karpacz Winter School of Theoretical Physics} (Acta Universitatis Wratislavensis no.~38, 1975); B.~S.~De Witt, in \textit{Quantum Gravity II}, eds. C.~Isham, R.~Penrose, and D.~Sciama (Oxford University Press, Oxford, 1982); L.~F.~Abbott, \textit{Nucl. Phys.} B\textbf{185}, 189 (1981).
\end{enumerate}
